\section{Assembly of the rational character sectors}

\begin{corollary}[AC-T46: the rational trace on the explicit prolate completion]
\label{AC-T46}
Use \(C_{\mathbb Q}\simeq
\widehat{\mathbb Z}^{\times}\times\mathbb R_+^\times\)
and the character convention of AC-T01.
For each nonprincipal primitive Dirichlet character
\(\chi\), use the target Hilbert space and the
full parity-selected prolate projection
\(Q_{\chi,\Lambda}\) of AC-T41, with
raw radii \((\Lambda,\Lambda/q_\chi)\).
For the principal character use the full completion
and projection \(P_\Lambda\) of AC-T45.
Their Hilbert direct sum has a bounded orthogonal
projection \(\mathbb Q_\Lambda\).
Let \(\mathbb R\) be the direct-sum centered action:
its nonprincipal summands have action
\(\overline{\chi(u)}V(r)\), and its principal
summand is (PC.3).

For \(h\in\mathcal D(C_{\mathbb Q})\), put
\[
 h_\chi(r)=\int_{\widehat{\mathbb Z}^{\times}}
                  h(u,r)\overline{\chi(u)}\,du .
 \tag{RA.1}
\]
Then \(\mathbb Q_\Lambda\mathbb R(h)\) has finite
rank.  Under \(\mathsf A\),
\[
 \begin{split}
 \operatorname{Tr}(\mathbb Q_\Lambda\mathbb R(h))
 ={}&2\log\Lambda\,h(1,1)
    -\sum_{\chi,\gamma}m_{\chi,\gamma}
               \int_0^\infty h_\chi(r)r^{i\gamma}\,\frac{dr}{r}\\
   &+\int_0^\infty h_1(r)
                    (r^{-1/2}+r^{1/2})\,\frac{dr}{r}
     +o(1).
 \end{split}
 \tag{RA.2}
\]
The character sum is finite for each fixed test.
The zero sums are absolutely convergent.
Equation (RA.2), followed by the explicit formula
in AC-T02, gives the geometric Weil distribution
with the same normalization.
This statement includes every rational unitary
norm twist and introduces no uniformity assertion
as the conductor varies.
\end{corollary}
\begin{proof}
The compact-unit part of a Bruhat test factors
through a finite quotient: each of the finitely
many locally constant factors in the definition
of \(\mathcal D\) has an open kernel of invariance,
and the intersection of those kernels is still
open.  Finite Fourier inversion on that quotient
therefore gives
\[
            h(u,r)=\sum_\chi h_\chi(r)\chi(u)
\]
with finitely many continuous compact-unit
characters.  Each such character is counted
once with its primitive conductor, as in AC-T01.
The integrated action is zero on all other
summands.  On each of its finitely many
nonzero summands the chosen projection has
finite rank.  This proves the trace assertion,
even though \(\mathbb Q_\Lambda\) itself need
not have finite rank.

Apply AC-T41 to each nonprincipal coefficient
and AC-T44--AC-T45 to the principal coefficient.
There are finitely many coefficients, so their
individual \(o(1)\) errors sum to \(o(1)\);
no interchange of an infinite conductor limit
is involved.  Finite Fourier inversion at
\(u=1,r=1\) gives
\(\sum_\chi h_\chi(1)=h(1,1)\).
This proves (RA.2).
The zero-count estimates and Schwartz decay
of each Mellin transform give the stated
absolute convergence.
A fixed unitary norm twist replaces a coefficient
by \(r^{it}h_\chi(r)\), still a smooth compact
norm test.  Thus the same proof covers it.
The equality with the centered geometric
distribution is the previously stated exact
Weil identity, not an additional cutoff limit.
\end{proof}

\begin{remark}[Exact scope of the assembled family]
The definition above gives one explicit rational
prolate completion, including its two principal
polar coordinates, and proves its actual trace.
In C99, VIII.(31)--(33), the displayed real
quantization and tensor map explicitly treat
the zeta sector.  The subsequent AC-T47 proves
the exact quotient metric and minimal-conductor-set
pair for the ramified sector construction.
AC-T48 proves equality of tested trace limits
with the angle cutoff for each fixed semilocal set.
Those are additional operator arguments; finite
Fourier assembly and the explicit formula alone
do not supply them.
\end{remark}
